\documentclass[12pt]{article}
\usepackage[utf8]{inputenc}
\usepackage[T1]{fontenc}
\usepackage[margin=1in]{geometry}
\usepackage{setspace}
\usepackage{xcolor}
\usepackage[colorlinks=true, urlcolor=blue!60!black, linkcolor=blue!60!black]{hyperref}
\setlength{\parindent}{0em}
\setlength{\parskip}{0.6em}
\renewcommand{\baselinestretch}{1.0}
\begin{document}
\begin{center}
{\LARGE \textit{Baker v.\ Carr}, 369 U.S.\ 186 (1962)\par}
\vspace{0.5em}
{\large Case Brief --- Two-Page Edition\par}
\end{center}

\section{Memory Jogger}
Equal Protection challenges to malapportioned state legislatures are justiciable; the case supplies the six-factor political question test (\href{https://supreme.justia.com/cases/federal/us/369/186/}{Justia}).

\section{Detailed Case Facts}
Tennessee had not reapportioned its legislature since 1901 despite a constitutional duty to do so every ten years. As population urbanized, about 37\% of voters elected 20 of 33 Senators, and Moore County (2,340 voters) had the same representation as Rutherford County (25,316). Urban voters led by Charles W. Baker sued under 42 U.S.C. \S~1983, alleging a ``debasement of their votes'' in violation of Equal Protection. Tennessee offered no initiative or referendum, and its courts had refused relief.

\section{Procedural History}
A three-judge district court dismissed under \textit{Colegrove v. Green} (1946), though it found a clear violation of the state constitution. On direct appeal and reargument, the Supreme Court reversed and remanded.

\section{Judicial Votes}
Reversed, 6--2. Brennan wrote for Warren, Black, Douglas, Clark, and Stewart; Douglas, Clark, and Stewart also concurred separately. Frankfurter, joined by Harlan, dissented, and Harlan dissented separately. Whittaker did not participate.

\section{Holding}
Federal courts have jurisdiction over such claims, voters alleging arbitrary vote dilution have standing, and an Equal Protection malapportionment claim is not a political question. The Court left the merits and remedy to the district court.

\section{Analysis of Opinions}
\textbf{Majority:} The political question doctrine concerns the judiciary's relationship to the coordinate federal branches, not to the States. Brennan distilled six factors from past cases; Guaranty Clause claims fail only because they implicate those factors, while this claim invoked ``well developed and familiar'' Equal Protection standards. \textbf{Clark:} the apportionment was ``a crazy quilt without rational basis,'' and courts were voters' only remedy. \textbf{Stewart:} the Court did not require equal weighting of votes. \textbf{Frankfurter:} this was a ``Guarantee Clause claim masquerading under a different label,'' a ``mathematical quagmire'' without standards. \textbf{Harlan:} the Fourteenth Amendment does not require equal weighting of votes.

\section{Examples: Future Applications}
\textbf{Same-Side:} under \textit{Reynolds v. Sims}, Alabama's county-based scheme with a 41-to-1 disparity is justiciable and unconstitutional. \textbf{Opposite-Side:} a challenge to the Senate's impeachment-trial committee is nonjusticiable under Factors 1 and 2 (\textit{Nixon v. United States}). \textbf{Fence-Sitter:} members of Congress sue to halt airstrikes lasting past 60 days, pitting statutory interpretation against foreign-relations commitment.

\section{Critique}
Brennan's claim that manageable standards existed was largely ipse dixit, forcing \textit{Reynolds} to adopt ``one person, one vote.'' Process theorists such as John Hart Ely praise \textit{Baker} for unblocking channels of democratic change, yet its manageable-standards prong later shielded partisan gerrymandering in \textit{Rucho v. Common Cause}.

\section{Key Quotations}
\begin{quote}
``Prominent on the surface of any case held to involve a political question is found a textually demonstrable constitutional commitment of the issue to a coordinate political department; or a lack of judicially discoverable and manageable standards for resolving it; or the impossibility of deciding without an initial policy determination of a kind clearly for non judicial discretion; or the impossibility of a court's undertaking independent resolution without expressing lack of the respect due coordinate branches of government; or an unusual need for unquestioning adherence to a political decision already made; or the potentiality of embarrassment from multifarious pronouncements by various departments on one question.''
\end{quote}
\noindent --- Justice Brennan, for the Court, 369 U.S. at 217.

\end{document}
